\chapter{Fundamentacao Teórica}

\section{Postulados da Mecânica Quântica}

A física quântica surgiu em meio a descoberta de que tanto ondas eletromagnéticas \
quanto partículas, poderiam apresentar comportamentos ondulatóricos e corpusculares a depender do experimento. \
Além disso, a formulaçao clássica baseada nas leis de newton falhava na descrição de tais sistemas, o que \
levou a um grande esforço da comunidade científica por um formalismo mais acertivo.

A descrição inicial mais completa da mecânica quântica aborda uma descrição probabilística dos sistemas em pequena escala. \
Essa descrição é muito motivada por fatos experimetais e abordagens teóricas antecedentes, \
tais como a dualidade onda-partícula, a ideia de onda de matéria de de Broglie, e a quantização \
da energia de Planck. O sucesso empírico dessa teoria faz com que seja amplamente utilizada \
em pesquisas atuais, em especial na física da matéria condensada.

Os aspectos mais fundamentais da mecânica quântica podem ser resumidos em seus postulados.

\subsection{Primeiro Postulado}

Assim como na mecânica clássica os sistemas possuem um conjunto de variáveis que definem o estado de um sistema, \
como por exemplo a posição e o momento, na mecânica quântica assumimos que nosso sistema também pode assumir \
um conjunto de estados representados como vetores no espaço de Hilbert denotados por $|k\rangle$. Contudo, diferente da mecânica clássica, \
um sistema quântico pode estar em uma superposição de estados físicos. Isso é descrito matemáticamente como uma \
combinação linear no caso de estados discretos,
\begin{equation}
    |\psi\rangle=\sum_k c_k |k\rangle,
\end{equation}
como por exemplo o spin do eletron, ou, quando lidamos com estados contínuos, como a posição de uma partícula, \
a superposição se da por uma integral sobre todos os estados,
\begin{equation}
    |\psi \rangle = \int dk ~c(k) |k \rangle.
\end{equation}

\subsection{Segundo Postulado}

De fato, o estado de um sistema quantico ser uma combinação linear de seus possíveis estados \
não tem um significado físico evidente. Portanto, essa interpretação é postulada de modo que possamos compreender \
o que o formalismo matemático representa na realidade. O segundo postulado da mecânica quântica afirma que se os estados são ortonormais, isso é, 
\begin{equation}
    \langle i | j \rangle = \delta_{ij}
\end{equation}
no caso discreto, e
\begin{equation}
    \langle \xi | \xi' \rangle = \delta(\xi-\xi')
\end{equation}
no caso contínuo, a probabilidade \
de se encontrar um sistema quântico no estado $|k\rangle$ é $|\langle k | \psi \rangle|^2$, os coeficientes da combinação linear, \
de modo que a probabilidade de achar o sistem em qualquer estado é
\begin{equation}
    \langle \psi | \psi \rangle = (\sum_k c_k^* \langle k |)(\sum_k c_k |k\rangle) = \sum_k |c_k|^2 = \sum_k |\langle k | \psi \rangle|^2 = 1
\end{equation}
ou
\begin{equation}
    \langle \psi | \psi \rangle = (\int dk ~c^*(k)\langle k |)(\int dk' ~c(k')|k'\rangle)=\int\int dk dk' ~c^*(k)c(k')\delta(k-k') = \int dk ~|c(k)|^2 = 1,
\end{equation}
como deve ser.

\subsection{Terceiro Postulado}

Uma medida de uma variável física como a posição, momento ou energia, em um sistema quântico, é representado por uma operação no estado do sistema. \
Essa operação é mediada por um objeto chamado operador, $\hat A$, que é linear e atua sobre o ket, de modo que os únicos valores possíveis de serem obtidos \
pela medição, $a_k$, são os que satisfazem a equação
\begin{equation}
    \hat A |\psi_k \rangle = a_k |\psi_k \rangle.
\end{equation}
Os estados $|\psi_k\rangle$ são conhecidos como autoestados de $\hat A$, enquanto que $a_k$ são seus autovalores.

\subsection{Quarto Postulado}

Os operadores $\hat x$ e $\hat p$ tem sua atuação conhecida na base de posições, 
\begin{equation}
    \hat x |\psi\rangle =  \int dx ~\psi(x) \hat x|x \rangle = \int dx ~\psi(x) x|x \rangle
\end{equation}
\begin{equation}
    \langle x|\hat x |\psi\rangle =  x \psi(x).
\end{equation}
Além disso, pode-se obter apartir dessa relação que
\begin{equation}
    \langle x|\hat p |\psi\rangle =  -i\hbar \frac{\partial}{\partial x} \psi(x).
\end{equation}
O quarto postulado afirma que se conhecemos uma função clássica $f(x,p)$ que depende da posição e do momento, e que queremos medi-la, podemos obter \
o seu respectivo operador na mecânica quântica promovendo a posição e o momento aos operadores de posição e moemnto,
\begin{equation}
    x, ~p \to \hat x, ~\hat p
\end{equation}
\begin{equation}
    f(x,p) \to f(\hat x,\hat p)
\end{equation}

\subsection{Quinto Postulado}

Se soubermos o estado de um sistema quântico em um dado instante, e queremos saber qual será seu estado \
em um período de tempo depois, recorremos ao quinto postulado que afirma que a evolução temporal de um sistema quântico \
é descrita pela equação de Schrodinger,
\begin{equation}
i\hbar \frac{\partial}{\partial t}|\psi\rangle = \hat H |\psi\rangle
\end{equation}
em que $\hat H$ é o operador Hamiltoniano,
\begin{equation}
    \hat H = \frac{\hat p^2}{2m} + \hat V,
\end{equation}
cujos autovalores são as possíveis energias do sistema.

\section{Estados Estacionários}

Em muitas situações importantes na física, o a energia potêncial não depende do tempo, \
ou simplesmente tratamos como se nao dependesse. Nesse caso especial, podemos reescrever \
equação de Schrodinger de forma mais conveniente.

Partindo da equação de Schrodinger e projetando na base de posição, temos que
\begin{equation}
    i\hbar \langle x | \frac{\partial}{\partial t}|\Psi\rangle = \langle x |(\frac{\hat p^2}{2m} + \hat V) |\Psi\rangle,
\end{equation}
logo
\begin{equation}
    i\hbar \frac{\partial}{\partial t}\Psi(x,t) = (- \frac{\hbar^2}{2m}\frac{\partial^2}{\partial x^2} + V(x)) \Psi(x,t)
\end{equation}
Assumindo uma solução do tipo $\Psi(x,t)=\psi(x)\phi(t)$, é possível separar em duas equações
\begin{equation}
\frac{i\hbar}{\phi(t)} \frac{\partial}{\partial t}\phi(t) = \frac{1}{\psi(x)}(- \frac{\hbar^2}{2m}\frac{\partial^2}{\partial x^2} + V(x)) \psi(x)=E.
\end{equation}
A equação para $\phi(t)$ pode ser resolvida por integração, 
\begin{equation}
    \phi(t) = e^{-i\frac{E}{\hbar}}
\end{equation}
enquanto que a equação para $\psi(x)$,
\begin{equation}
    (- \frac{\hbar^2}{2m}\frac{\partial^2}{\partial x^2} + V(x)) \psi(x) = E \psi(x),
\end{equation}
depende do potêncial, e é conhecida como equação de Schrodinger independente do tempo. \
Ainda de maneira mais geral podemos escrever essa equação como
\begin{equation}
    (- \frac{\hbar^2 \hat p^2}{2m} + V(\hat x)) |\psi\rangle = E |\psi\rangle,
\end{equation}
com $|\Psi(t)\rangle = e^{-i\frac{E}{\hbar}} |\psi\rangle$. \
Esse tipo de equação é chamado de solução estacionária, e possui algumas características importantes.



